\section{Colored covers and exact packet reconstruction}
\label{sec:cover-reconstruction}

We prepare the passage from the boundary slots supplied by
\Cref{thm:torus-compression} to a decrease of covering order.
The reconstruction must retain the original open-cover information
exactly.  A coloring of a subdivided nerve makes this possible: a
boundary output leaves at least one color absent from its packet.
To preserve the original zero coordinates, reconstruction discards
packets of small total weight.  On every retained packet, fixing the
original boundary value prevents an originally absent color from
appearing.

\subsection{Polyhedral models for finite covers}

For a finite simplicial complex $S$, write $\lambda_v:S\to[0,1]$
for the barycentric coordinate of a vertex $v$.  The \emph{carrier}
of $y\in S$ is the unique simplex whose relative interior contains
$y$; its vertices are precisely those for which $\lambda_v(y)>0$.
In particular, its dimension is the number of positive vertex
coordinates minus one.

We will use the following elementary refinement principle.

\Needspace{6\baselineskip}
\begin{lemma}[Polyhedral refinement]
\label{lem:polyhedral-refinement}
Let $Y$ be a compact space, let $K$ be a finite polyhedron of
dimension at most $q$, and let $\Phi:Y\to K$ be continuous.
Let $\mathcal V$ be a finite open cover of $Y$.  Suppose that every
point of $\Phi(Y)$ has an open neighborhood $W$ in $K$ such that
$\Phi^{-1}(W)$ is contained in a member of $\mathcal V$.
Then $D(\mathcal V)\le q$.
\end{lemma}

\begin{proof}
Choose finitely many such neighborhoods covering the compact set
$\Phi(Y)$, and add $K\setminus\Phi(Y)$ to obtain a finite open
cover of $K$.  A sufficiently fine simplicial subdivision has its
open vertex stars subordinate to this cover.  These stars have order
at most $q$: a collection of stars with nonempty intersection
corresponds to vertices of a single simplex.  Their pullbacks under
$\Phi$, with empty sets omitted, form a finite open refinement of
$\mathcal V$ of order at most $q$.
\end{proof}

\begin{lemma}[Colored subordinate nerve]
\label{lem:colored-subordinate-nerve}
Let $\mathcal V$ be a finite open cover of a compact metrizable
space $Y$, and put $m=D(\mathcal V)$.  There are a finite simplicial
complex $S$ of dimension at most $m$, a continuous map
$\phi:Y\to S$, a coloring of the vertices of $S$ by
$1,\ldots,m+1$ that is injective on every simplex, and, for each
vertex $v$, a member $V(v)\in\mathcal V$ such that
\begin{equation}
\label{eq:vertex-subordination}
 \{y\in Y:\lambda_v(\phi(y))>0\}\subset V(v).
\end{equation}
\end{lemma}

\begin{proof}
Choose a finite open refinement of $\mathcal V$ of order $m$ and
a subordinate partition of unity.  The partition defines a map
to the nerve $S_0$ of that refinement, whose dimension is at most
$m$.  Assign each vertex of $S_0$ a containing member of
$\mathcal V$.

Take $S$ to be the barycentric subdivision of $S_0$, and regard
the same map as a map to $S$ under the canonical identification
of their realizations.  A vertex of $S$ is a nonempty face $F$ of
$S_0$; give it color $\dim F+1$ and choose one original vertex
$v_F\in F$.  Faces in a simplex of $S$ form a strictly increasing
chain, so their colors are distinct.  If the coordinate at the
barycenter of $F$ is positive, then the original coordinate at
$v_F$ is positive: the contribution of that barycenter to this
coordinate is its positive coefficient divided by the number of
vertices of $F$.  Assign to the new vertex the member of
$\mathcal V$ already assigned to $v_F$.  This proves
\Cref{eq:vertex-subordination}.
\end{proof}

For completeness, these models also explain the finite-cover
limits used below.  Given two finite covers, take subordinate
nerve maps and their product.  Around every point of the product,
choose a positive vertex in each factor.  The corresponding product
of positive-coordinate neighborhoods pulls back into a member of
the joined cover.  By \Cref{lem:polyhedral-refinement},
\begin{equation}
\label{eq:cover-subadditivity}
 D(\mathcal V\vee\mathcal W)
 \le D(\mathcal V)+D(\mathcal W).
\end{equation}
Here a product of finite polyhedra has dimension at most the sum
of their dimensions, as is also seen from its product cells.
For a finite open cover $\mathcal U$ of $X$, write
\[
 \mathcal U^{[a,b)}=\bigvee_{s=a}^{b-1}h^{-s}\mathcal U,
 \qquad m_k=D(\mathcal U^{[0,k)}).
\]
Invariance of $D$ under homeomorphisms and
\Cref{eq:cover-subadditivity} give $m_{k+l}\le m_k+m_l$.
Consequently
\begin{equation}
\label{eq:finite-cover-rate}
 d_{\mathcal U}
   =\lim_{k\to\infty}\frac{m_k}{k}
   =\inf_{k\ge1}\frac{m_k}{k}
   \le D(\mathcal U)<\infty.
\end{equation}
This finiteness concerns one cover and does not require finite
mean dimension or finite covering dimension of $X$.

\subsection{Color packets and exact reconstruction}

Recall $C_t=[-1,1]^t$ and fix a homeomorphism
$\kappa:\Delta^t\to C_t$ carrying the whole boundary of the
simplex onto the whole boundary of the cube.  Such a map may, for
example, be obtained by radial extension from interior points of
these two convex bodies.

\begin{lemma}[Packet reconstruction]
\label{lem:packet-reconstruction}
Let $S$ be a finite simplicial complex of dimension at most $m$
whose vertices are colored by $1,\ldots,m+1$, with distinct colors
on each simplex.  Fix an integer $t\ge1$ and put
\[
 g=\left\lfloor\frac{m+1}{t+1}\right\rfloor.
\]
There exist continuous maps $\mathcal F_j:S\to C_t$,
$1\le j\le g$, with the following property.

Let $Y$ be any topological space, let $y:Y\to S$ be continuous,
and let $v_j:Y\to C_t$ be continuous replacements satisfying
\begin{equation}
\label{eq:packet-boundary-condition}
 v_j(z)=\mathcal F_j(y(z))
 \quad\text{whenever }\mathcal F_j(y(z))\in\partial C_t.
\end{equation}
There is then a continuous reconstructed map
$\widetilde y:Y\to S$ such that, for every $z\in Y$,
\begin{enumerate}
 \item $\lambda_v(\widetilde y(z))>0$ implies
       $\lambda_v(y(z))>0$ for every vertex $v$;
 \item if $s(z)=\#\{j:v_j(z)\in\partial C_t\}$, then the carrier
       of $\widetilde y(z)$ has dimension at most $m-s(z)$.
\end{enumerate}
In particular, reconstruction preserves every vertex
subordination relation such as \Cref{eq:vertex-subordination}.
\end{lemma}

\begin{proof}
For a color $\ell$, let
\[
 c_\ell(y)=\sum_{v\text{ of color }\ell}\lambda_v(y).
\]
These continuous color weights sum to one.  At any point of $S$
at most one vertex of a given color has positive weight.
Partition the first $(t+1)g$ colors into packets
\[
 I_j=\{(j-1)(t+1)+1,\ldots,j(t+1)\},
 \qquad 1\le j\le g,
\]
and retain the remaining colors individually.  Let $w_j(y)$ be
the vector of the $t+1$ color weights in $I_j$, and let
$M_j(y)$ be their sum.  At positive mass, the proportions $w_j/M_j$
record each missing color by a zero coordinate, but division by
$M_j$ is undefined at zero mass.  We first regularize these
proportions near zero to obtain continuous slot maps.  During
reconstruction we discard packets in this range, so any colors
introduced by regularization do not enter the reconstructed point.
Choose
\begin{equation}
\label{eq:packet-threshold}
 0<\tau<\frac{1}{4(m+1)},
 \qquad b=\left(\frac{1}{t+1},\ldots,\frac{1}{t+1}\right)
 \in\Delta^t.
\end{equation}
Define a continuous simplex-valued function by
\begin{equation}
\label{eq:regularized-proportions}
 q_j(y)=
 \begin{cases}
  w_j(y)/M_j(y),&M_j(y)\ge\tau,\\[2mm]
  b+\bigl(w_j(y)-M_j(y)b\bigr)/\tau,&M_j(y)\le\tau.
 \end{cases}
\end{equation}
The formulas agree at $M_j=\tau$.  For $0<M_j\le\tau$,
the second is the convex combination of $b$ and $w_j/M_j$
with coefficients $1-M_j/\tau$ and $M_j/\tau$; at $M_j=0$
it equals $b$.  Thus it takes values in $\Delta^t$ throughout.
Set $\mathcal F_j=\kappa\circ q_j$.

To implement this deletion, choose a continuous function
$\eta:[0,1]\to[0,1]$ that is zero
on $[0,\tau]$ and one on $[2\tau,1]$.  For data as in the
statement, give packet $j$ the unnormalized new color vector
\begin{equation}
\label{eq:reconstructed-packet}
 \widehat w_j(z)
   =\eta(M_j(y(z)))M_j(y(z))\kappa^{-1}(v_j(z)),
\end{equation}
and give each leftover color its original weight.  Denote all
these new color weights by $\widehat c_\ell(z)$, and their sum by
$Z(z)$.  Then
\[
 Z(z)=1-\sum_{j=1}^g
       (1-\eta(M_j(y(z))))M_j(y(z))
       \ge1-2g\tau>0.
\]
The last inequality follows from \Cref{eq:packet-threshold}.
Thus division by $Z$ is defined and continuous, even if some
packets have been discarded completely.

We first verify that a missing original color stays missing.
For a leftover color this is immediate.  For a packet whose
factor $\eta(M_j)M_j$ is zero it is also immediate.  In the
remaining case $M_j>\tau$, so $q_j=w_j/M_j$ exactly.  If one
original color in that packet has weight zero, $q_j$ lies on
$\partial\Delta^t$ and $\mathcal F_j$ lies on $\partial C_t$.
\Cref{eq:packet-boundary-condition} forces
$\kappa^{-1}(v_j)=q_j$, and the new weight of that color is
zero.  We have proved
\begin{equation}
\label{eq:no-new-colors}
 c_\ell(y(z))=0\quad\Longrightarrow\quad
 \widehat c_\ell(z)=0.
\end{equation}

Assign the normalized color weights to the originally present
vertices.  More explicitly, for a vertex $v$ of color $\ell$,
put
\begin{equation}
\label{eq:reconstructed-vertices}
 \widetilde\lambda_v(z)=
 \begin{cases}
 \displaystyle
 \frac{\lambda_v(y(z))}{c_\ell(y(z))}
 \frac{\widehat c_\ell(z)}{Z(z)},
       &c_\ell(y(z))>0,\\[3mm]
 0,    &c_\ell(y(z))=0.
 \end{cases}
\end{equation}
This function is continuous at points where $c_\ell=0$ as well:
the first ratio lies in $[0,1]$, while the second tends to zero
by \Cref{eq:no-new-colors} and the positive lower bound for $Z$.
The new vertex weights are nonnegative, sum to one, and have
support contained in the original carrier.  They therefore
define a continuous map $\widetilde y$ into $S$, with the first
asserted property.

Finally, if $v_j(z)\in\partial C_t$, the vector
$\kappa^{-1}(v_j(z))$ has a zero coordinate.  Thus packet $j$
contributes at most $t$ positive colors in
\Cref{eq:reconstructed-packet}; if the packet is discarded,
it contributes none.  The $s(z)$ successful packets use disjoint
sets of colors.  Hence at most $m+1-s(z)$ colors, and therefore
at most $m+1-s(z)$ vertices, are positive in the reconstructed
point.  Its carrier has dimension at most $m-s(z)$.
\end{proof}

Notice that the loss in \Cref{lem:packet-reconstruction} is measured
from the common bound $m$, not from the dimension of the original
carrier.  A packet already missing a color may already have a
boundary input, and is counted correctly by this formulation.
Nor do discarded packets cause a loss of the asserted gain: a
discarded packet has all of its colors absent.
